% Step 1 guide: Mileti's formulas in Lean 4 (beginner edition)
% Build:  cd docs/guides && lualatex step1-guide.tex && lualatex step1-guide.tex
\documentclass[11pt,a4paper]{article}
\usepackage[margin=2.3cm]{geometry}
\usepackage{fontspec}
\directlua{luaotfload.add_fallback("fb",{"STIXTwoMath-Regular.otf:mode=node;"})}
\setmainfont{DejaVuSerif.ttf}[RawFeature={fallback=fb},BoldFont=DejaVuSerif-Bold.ttf,ItalicFont=DejaVuSerif-Italic.ttf,Scale=0.92]
\setsansfont{DejaVuSans.ttf}[RawFeature={fallback=fb},BoldFont=DejaVuSans-Bold.ttf]
\setmonofont{DejaVuSansMono.ttf}[RawFeature={fallback=fb},BoldFont=DejaVuSansMono-Bold.ttf,Scale=0.82]
\usepackage{xcolor}
\usepackage{fvextra}
\usepackage[most]{tcolorbox}
\usepackage{booktabs,tabularx,array}
\usepackage{enumitem}
\usepackage{titlesec}
\usepackage{fancyhdr}
\usepackage[hidelinks]{hyperref}
\usepackage{parskip}

\definecolor{accent}{HTML}{1F5FA8}
\definecolor{codebg}{HTML}{F5F7FA}
\definecolor{outbg}{HTML}{EEF6EE}
\definecolor{tipc}{HTML}{1F5FA8}
\definecolor{okc}{HTML}{2E7D32}
\definecolor{warnc}{HTML}{B26A00}
\definecolor{tryc}{HTML}{6A1B9A}

\titleformat{\section}{\sffamily\Large\bfseries\color{accent}}{Part \thesection}{0.8em}{}
\titleformat{\subsection}{\sffamily\large\bfseries}{\thesubsection}{0.6em}{}
\titleformat{\subsubsection}{\sffamily\normalsize\bfseries}{}{0em}{}
\setcounter{section}{-1}

\pagestyle{fancy}
\fancyhf{}
\fancyhead[L]{\sffamily\small Step 1: Mileti's formulas in Lean}
\fancyhead[R]{\sffamily\small mileti-kg-prover}
\fancyfoot[C]{\sffamily\small\thepage}
\renewcommand{\headrulewidth}{0.3pt}

\DefineVerbatimEnvironment{lean}{Verbatim}{fontsize=\small,breaklines,frame=leftline,framerule=2pt,rulecolor=\color{accent},xleftmargin=4pt,framesep=8pt}
\DefineVerbatimEnvironment{shell}{Verbatim}{fontsize=\small,breaklines,frame=leftline,framerule=2pt,rulecolor=\color{gray},xleftmargin=4pt,framesep=8pt}
\DefineVerbatimEnvironment{output}{Verbatim}{fontsize=\small,breaklines,frame=leftline,framerule=2pt,rulecolor=\color{okc},xleftmargin=4pt,framesep=8pt}
\newcommand{\code}[1]{\texttt{#1}}

\newtcolorbox{tip}[1][Tip]{enhanced,breakable,colback=tipc!5,colframe=tipc,fonttitle=\sffamily\bfseries,title={#1},boxrule=0.6pt,arc=2pt}
\newtcolorbox{okbox}[1][Check]{enhanced,breakable,colback=okc!5,colframe=okc,fonttitle=\sffamily\bfseries,title={✓\ #1},boxrule=0.6pt,arc=2pt}
\newtcolorbox{warn}[1][Watch out]{enhanced,breakable,colback=warnc!6,colframe=warnc,fonttitle=\sffamily\bfseries,title={#1},boxrule=0.6pt,arc=2pt}
\newtcolorbox{try}[1][Try it yourself]{enhanced,breakable,colback=tryc!5,colframe=tryc,fonttitle=\sffamily\bfseries,title={#1},boxrule=0.6pt,arc=2pt}

\setlist{itemsep=2pt,topsep=4pt}

\begin{document}

\begin{titlepage}
\vspace*{3cm}
{\sffamily\color{accent}\Huge\bfseries Step 1: Mileti's formulas in Lean 4\par}
\vspace{0.6cm}
{\sffamily\Large A beginner's guide to the first part of the project\par}
\vspace{1.2cm}
{\large Project: \emph{Legal moves and book knowledge: can LLMs build textbook deductions that Lean accepts?}\par}
\vspace{0.4cm}
{\large Repository: \code{github.com/ZaDrMeister/mileti-kg-prover}\par}
\vspace{2cm}
\begin{tcolorbox}[colback=accent!4,colframe=accent,boxrule=0.6pt,arc=2pt]
\textbf{What you'll have at the end}
\begin{itemize}
\item \code{lean/MiletiLogic/Syntax.lean}, written by you: Mileti's propositional formulas as a Lean type, with readable notation and two recursive functions, all checked by Lean.
\item Your first commits on GitHub.
\item Your first research-diary entries.
\end{itemize}
\textbf{Time:} about 6--8 hours, split over three days (see the plan on page \pageref{plan}).\\[4pt]
\textbf{Assumes:} no Lean experience, basic programming (you've used Python or Java), and Mileti Ch.~2--3 at hand.
\end{tcolorbox}
\vfill
{\small Written 2026-11-25. Every Lean snippet in this guide was compiled with Lean v4.28.0, the version in your project. If a snippet fails for you, look for a typo or a setup problem first.\par}
\end{titlepage}

\tableofcontents
\newpage

% ------------------------------------------------------------------
\section{Before you start}

\subsection{The one rule}
\textbf{You write everything in \code{Syntax.lean}.} This guide teaches with \emph{other} examples (colors, shapes, arithmetic expressions) that have the same structure as formulas. Your job is to translate them. That's on purpose:
\begin{itemize}
\item The project's main claim is that \emph{you} formalized Mileti's system. PhD interviewers will ask you to explain your code line by line.
\item Translating an example is where the learning happens. Copying doesn't teach you anything.
\end{itemize}
When you're stuck, ask Claude \emph{what an error means} or \emph{for a hint}. Don't ask for the code.

\subsection{How to read this guide}
Each part uses these boxes:

\begin{tip}[Tip] Extra explanation or a shortcut. \end{tip}
\begin{try} A small exercise in your scratch file. Do it: each one prepares a piece of your real file. \end{try}
\begin{okbox}[Check] How to know a step worked. Don't go on until it passes. \end{okbox}
\begin{warn} A common mistake. \end{warn}

Code you type looks like this (blue bar):
\begin{lean}
#eval 2 + 3
\end{lean}
What Lean shows back looks like this (green bar):
\begin{output}
5
\end{output}
Terminal commands have a grey bar:
\begin{shell}
cd ~/Desktop/grad-school-project/mileti-kg-prover
\end{shell}

\subsection{The three-day plan}\label{plan}
\begin{center}
\begin{tabularx}{\linewidth}{@{}l>{\raggedright\arraybackslash}Xl@{}}
\toprule
\textbf{Day} & \textbf{Parts} & \textbf{Time}\\
\midrule
Day 1 & Part 1 (why) + Part 2 (setup) + Lessons A--C in Part 3 & 2--2.5 h\\
Day 2 & Lessons D--F in Part 3 + Part 4 (reading Mileti) & 2--2.5 h\\
Day 3 & Part 5 (build \code{Syntax.lean}) + Part 6 (finish and commit) & 2--3 h\\
\bottomrule
\end{tabularx}
\end{center}
End every day with a diary entry (Part 8). It takes 10 minutes.

% ------------------------------------------------------------------
\section{What you're building, and why}

\subsection{Formulas in Mileti}
In Mileti (Ch.~2 and §3.1), a \textbf{formula} of propositional logic is built from:
\begin{itemize}
\item \textbf{atoms}: propositional symbols such as $A$, $B$, $C$ (the book uses a set $P$ of them);
\item \textbf{connectives} that build bigger formulas from smaller ones: $\neg\varphi$, $\varphi\wedge\psi$, $\varphi\vee\psi$, $\varphi\to\psi$.
\end{itemize}
Mileti defines the set of formulas with a \textbf{generating system}: start with the atoms, keep applying the connectives, and take everything you can reach. He writes formulas as \emph{strings} in Polish (prefix) notation, for example \code{∧AB} for $A\wedge B$. Because strings could in principle be read in more than one way, the book has to \emph{prove} \textbf{unique readability}: every formula is built in exactly one way.

\subsection{Formulas in Lean: trees}
In Lean you don't use strings. You describe formulas as \textbf{trees}:
\begin{center}
\begin{BVerbatim}
          →            this tree is  (A ∧ B) → (A ∨ B)
        /   \
       ∧     ∨         the top node is the "main connective"
      / \   / \
     A   B A   B       the leaves are atoms
\end{BVerbatim}
\end{center}
Lean's \code{inductive} keyword creates exactly this kind of tree type. Lean \emph{guarantees} two things about it:
\begin{enumerate}
\item Trees built with different constructors are different (an $\wedge$-tree is never an $\vee$-tree).
\item Each constructor is one-to-one: if $\varphi_1\wedge\psi_1$ and $\varphi_2\wedge\psi_2$ are the same tree, then $\varphi_1=\varphi_2$ and $\psi_1=\psi_2$.
\end{enumerate}
Those are exactly Mileti's \textbf{freeness} conditions for a generating system. So with an inductive type, \emph{unique readability holds automatically}. That's the first insight of the project, and you'll write it up in your own words in Part~4.

\subsection{Where step 1 fits}
\begin{center}
\begin{tabularx}{\linewidth}{@{}l>{\raggedright\arraybackslash}X@{}}
\toprule
\textbf{Step} & \textbf{What}\\
\midrule
\textbf{1 (this guide)} & \code{Formula}: what a formula \emph{is}\\
2 & \code{Derivable}: Mileti's 14 deduction rules (§3.4), built on \code{Formula}\\
3 & Python gates: one per rule, the legal-move check\\
later & the LLM loop, the Lean check of every proof, experiments (conditions A, B, D, E, W)\\
\bottomrule
\end{tabularx}
\end{center}
Everything later depends on \code{Formula}, so it's worth getting it right.

% ------------------------------------------------------------------
\section{Setup (30 minutes)}

\subsection{Open Terminal and check that Lean works}
Open the \textbf{Terminal} app (⌘+Space, type ``Terminal'', Enter). Type each line and press Enter:
\begin{shell}
cd ~/Desktop/grad-school-project/mileti-kg-prover/lean
~/.elan/bin/lake build
\end{shell}
\code{cd} means ``change directory'' (go into that folder). \code{lake} is Lean's build tool: it checks every Lean file in the project.
\begin{okbox}
You see \code{Build completed successfully}. The first run can take a minute.
\end{okbox}
\begin{warn}[If it fails]
Copy the \emph{whole} output and ask Claude what it means. Don't continue until the build works.
\end{warn}

\subsection{Install the Lean extension in VS Code}
\begin{enumerate}
\item Open VS Code.
\item Open the Extensions panel: ⇧⌘X (Shift+Command+X).
\item Search \textbf{lean4}. Install the one published by \textbf{leanprover}.
\item Restart VS Code if it asks you to.
\end{enumerate}

\subsection{Open the right folder}
\textbf{File → Open Folder…} and choose
\code{\textasciitilde/Desktop/grad-school-project/mileti-kg-prover/lean}.
\begin{warn}
Open the \code{lean} folder itself, not the whole project. The Lean extension looks for \code{lakefile.lean} in the folder you open. With the wrong folder, Lean can't find your files and the infoview stays empty.
\end{warn}

\subsection{The infoview: where Lean talks to you}
Open \code{MiletiLogic/Syntax.lean} from the sidebar. Then open the infoview: press ⇧⌘P, type \code{Lean 4: Infoview: Toggle}, and press Enter. A panel opens on the right.

The infoview shows, for the line your cursor is on:
\begin{itemize}
\item results of \code{\#eval} and \code{\#check};
\item \textbf{errors} (also underlined in red in the editor);
\item later (step 2), the \emph{goal} you're proving.
\end{itemize}
Errors also appear in VS Code's \textbf{Problems} panel (⇧⌘M).

\subsection{Make a scratch file}
Create a new file and save it as \code{lean/Scratch.lean}. That's where you'll practise. It's already listed in \code{.gitignore}, so git ignores it. Type:
\begin{lean}
#eval 2 + 3
\end{lean}
Put the cursor on that line.
\begin{okbox}
The infoview shows \code{5}. Your setup works.
\end{okbox}

\subsection{Typing logic symbols}
The Lean extension turns backslash codes into symbols. Type the code, then a space:
\begin{center}
\begin{tabular}{@{}llll@{}}
\toprule
\textbf{Type} & \textbf{Get} & \textbf{Type} & \textbf{Get}\\
\midrule
\code{\textbackslash to} & → & \code{\textbackslash ne} & ≠\\
\code{\textbackslash and} & ∧ & \code{\textbackslash phi} & φ\\
\code{\textbackslash or} & ∨ & \code{\textbackslash psi} & ψ\\
\code{\textbackslash not} & ¬ & \code{\textbackslash sim} & ∼\\
\code{\textbackslash bigwedge} & ⋀ & \code{\textbackslash bigvee} & ⋁\\
\code{\textbackslash longrightarrow} & ⟶ & \code{\textbackslash vdash} & ⊢\\
\bottomrule
\end{tabular}
\end{center}
Hover over any symbol in the editor to see how to type it.

% ------------------------------------------------------------------
\section{A Lean crash course in six lessons}
Do every lesson in \code{Scratch.lean}. \textbf{Type the examples; don't paste them.} Typing is how the syntax sticks. After each block, put your cursor on the \code{\#eval} lines and compare the infoview with the green output shown here.

\subsection{Lesson A: values, functions, checks (20 min)}
\begin{lean}
#eval 2 + 3                        -- run code
#eval "hello " ++ "Lean"           -- ++ joins strings
#check Nat                         -- what type is this?
#check (2 + 2 = 4)                 -- a statement is also something with a type

def double (n : Nat) : Nat := n + n
#eval double 21
#check double

example : double 2 = 4 := by rfl
\end{lean}
\begin{output}
5
"hello Lean"
Nat : Type
2 + 2 = 4 : Prop
42
double (n : Nat) : Nat
\end{output}
\textbf{Line by line:}
\begin{itemize}
\item \code{--} starts a comment. Lean ignores the rest of the line.
\item \code{\#eval} \emph{runs} something and shows the result. \code{\#check} shows its \emph{type}. Commands starting with \code{\#} are for you; they don't become part of the program.
\item \code{Nat} is the type of natural numbers $0,1,2,\dots$ \code{Prop} is the type of \emph{statements}. In Lean, ``$2+2=4$'' is a thing of type \code{Prop}. That's how Lean will later talk about ``$S\vdash\varphi$ is derivable''.
\item \code{def double (n : Nat) : Nat := n + n} reads: ``define \code{double}; it takes an input \code{n} of type \code{Nat}, returns a \code{Nat}, and the result is \code{n + n}.'' The pattern is always \code{def name (input : Type) : OutputType := body}.
\item \code{example : double 2 = 4 := by rfl} asks Lean to \emph{check} a fact. \code{rfl} (reflexivity) means ``compute both sides; they're the same.'' If the fact is false, you get a red error and no output.
\end{itemize}
\begin{try}
\begin{enumerate}
\item Write \code{def triple (n : Nat) : Nat := …} and \code{example : triple 3 = 9 := by rfl}.
\item Change \code{9} to \code{10}. Read the error message, then change it back.
\end{enumerate}
\end{try}

\subsection{Lesson B: your first \texttt{inductive} type (20 min)}
\begin{lean}
inductive Color where
  | red
  | green
  | blue
  deriving Repr, DecidableEq

#eval Color.red
example : Color.red ≠ Color.blue := by decide

def Color.next : Color → Color
  | .red   => .green
  | .green => .blue
  | .blue  => .red

#eval Color.red.next
\end{lean}
\begin{output}
Color.red
Color.green
\end{output}
\textbf{Line by line:}
\begin{itemize}
\item \code{inductive Color where} creates a new type. Each line starting with \code{|} is a \textbf{constructor}: one of the \emph{only} ways to make a \code{Color}. Nothing else is a \code{Color}.
\item \code{deriving Repr} lets \code{\#eval} \emph{print} a value of this type.
\item \code{deriving DecidableEq} lets Lean \emph{decide} whether two values are equal. \code{by decide} uses it. \textbf{You'll need both for formulas}: the gates and Lean must be able to ask ``is this formula the same as that one?''
\item \code{Color → Color} is the type ``function from Color to Color''.
\item The function is defined by \textbf{pattern matching}: one line per constructor. \code{.red} is short for \code{Color.red}; Lean knows the input is a \code{Color}.
\item \code{Color.red.next} is the same as \code{Color.next Color.red}. Naming the function \code{Color.next} lets you use the dot form.
\end{itemize}
\begin{try}
Add a fourth constructor \code{| yellow} to \code{Color}. Lean now marks \code{Color.next} with \code{Missing cases}. Fix it so the colors cycle red → green → blue → yellow → red, and check that \code{\#eval Color.yellow.next} shows \code{Color.red}.
\end{try}
\begin{tip}[Why \texttt{Missing cases} is good news]
Lean refuses functions that forget a case. When you have a function over all formulas (or, in step 2, over all 14 rules), Lean will tell you if you forgot one.
\end{tip}

\subsection{Lesson C: constructors that carry data (15 min)}
\begin{lean}
inductive Shape where
  | circle : Nat → Shape          -- a circle needs one number (radius)
  | rect   : Nat → Nat → Shape    -- a rectangle needs two numbers
  deriving Repr, DecidableEq

#eval Shape.rect 2 3

def Shape.area : Shape → Nat
  | .circle r => 3 * r * r        -- name the data: r
  | .rect w h => w * h

#eval (Shape.rect 2 3).area
\end{lean}
\begin{output}
Shape.rect 2 3
6
\end{output}
\textbf{Line by line:}
\begin{itemize}
\item \code{circle : Nat → Shape} reads: ``give me a \code{Nat}, I give you a \code{Shape}.''
\item \code{rect : Nat → Nat → Shape} reads: ``give me two \code{Nat}s, I give you a \code{Shape}.'' The arrows chain: the last type is the result, the others are inputs.
\item In the match, \code{.circle r} \emph{names} the number inside the circle \code{r}, so the right side can use it.
\end{itemize}
\begin{tip}[Connection to Mileti]
An \textbf{atom} is like \code{circle}: a constructor carrying a \emph{name} (a number or a string), not another formula. In Part~5 you'll decide which.
\end{tip}

\subsection{Lesson D: lists (15 min)}
Your \code{subformulas} function will return a \emph{list} of formulas, so learn lists first.
\begin{lean}
#eval [1, 2] ++ [3]              -- join two lists
#eval 0 :: [1, 2]                -- put one item at the front
#eval [1, 2, 3].length
#eval [1, 1, 2].eraseDups        -- remove duplicates
#eval [1, 2, 3].contains 2
\end{lean}
\begin{output}
[1, 2, 3]
[0, 1, 2]
3
[1, 2]
true
\end{output}
\begin{itemize}
\item \code{[a, b, c]} is a list. \code{List Nat} is the type ``list of natural numbers''; \code{List Formula} will be ``list of formulas''.
\item \code{x :: xs} is ``x, then the list xs''. \code{xs ++ ys} joins two lists.
\item Lists keep \emph{order} and \emph{duplicates}. Sets don't. You'll decide which one matches Mileti's subformulas.
\end{itemize}

\subsection{Lesson E: a type that contains itself (a tree). The key lesson (45 min)}
This has exactly the shape you need for formulas, but for arithmetic expressions. \textbf{Spend the most time here.}
\begin{lean}
namespace Demo

inductive Expr where
  | num   : Nat → Expr               -- base case: like an atom
  | neg   : Expr → Expr              -- one smaller Expr inside: like ¬
  | plus  : Expr → Expr → Expr       -- two smaller Exprs: like ∧, ∨, →
  | times : Expr → Expr → Expr
  deriving DecidableEq, Repr

open Expr      -- now `num 1` works instead of `Expr.num 1`

def size : Expr → Nat               -- number of nodes in the tree
  | num _     => 1                   -- _ : "some data I don't need"
  | neg e     => size e + 1          -- calls itself on the smaller part
  | plus a b  => size a + size b + 1
  | times a b => size a + size b + 1

def height : Expr → Nat             -- longest path from the top down
  | num _     => 0
  | neg e     => height e + 1
  | plus a b  => max (height a) (height b) + 1
  | times a b => max (height a) (height b) + 1

def subexprs : Expr → List Expr     -- every piece, including the whole
  | num n     => [num n]
  | neg e     => neg e :: subexprs e
  | plus a b  => plus a b :: (subexprs a ++ subexprs b)
  | times a b => times a b :: (subexprs a ++ subexprs b)

#eval size (plus (num 1) (neg (num 2)))
#eval height (plus (num 1) (neg (num 2)))
#eval subexprs (plus (num 1) (num 1))
#eval (subexprs (plus (num 1) (num 1))).eraseDups

example : plus (num 1) (num 2) ≠ plus (num 2) (num 1) := by decide

end Demo
\end{lean}
\begin{output}
4
2
[Demo.Expr.plus (Demo.Expr.num 1) (Demo.Expr.num 1), Demo.Expr.num 1, Demo.Expr.num 1]
[Demo.Expr.plus (Demo.Expr.num 1) (Demo.Expr.num 1), Demo.Expr.num 1]
\end{output}
\textbf{What's new, carefully:}
\begin{itemize}
\item \code{neg : Expr → Expr}: a constructor can take \emph{another} \code{Expr}. That makes \code{Expr} a \textbf{tree}: \code{plus (num 1) (neg (num 2))} is the tree with $+$ on top, $1$ on the left and $-2$ on the right.
\item \textbf{Recursion.} \code{size} calls \code{size} on the smaller parts. Lean accepts this because the parts are \emph{structurally smaller} (they're inside the tree), so the recursion must end. This is the same as Mileti's ``definition by recursion on formulas'', and it's why Mileti needs freeness: in a free system, recursive definitions are well-defined.
\item \textbf{Work through \code{size} by hand once:}
\[\code{size}(\code{plus}(1, \code{neg}(2))) = \code{size}(1) + \code{size}(\code{neg}(2)) + 1 = 1 + (1+1) + 1 = 4.\]
\item \code{namespace Demo … end Demo} groups names: inside, \code{Expr} is really \code{Demo.Expr}. Your file uses \code{namespace MiletiLogic}.
\item The \code{example … ≠ … := by decide} shows that $1+2$ and $2+1$ are \emph{different trees}, even though they have the same value. Formulas are like this too: $A\wedge B$ and $B\wedge A$ are different formulas.
\end{itemize}
\begin{try}
\begin{enumerate}
\item Write \code{countNums : Expr → Nat}, which counts only the \code{num} leaves. Check \code{countNums (plus (num 1) (neg (num 2)))} gives \code{2}.
\item Compute \code{height} of \code{times (plus (num 1) (num 2)) (num 3)} \emph{on paper}, then check with \code{\#eval}.
\item Write \code{subexprs} of \code{neg (plus (num 1) (num 2))} on paper, in order, then check.
\end{enumerate}
\end{try}
\begin{tip}[{If you can do Lesson E, you can do step 1}]
\code{Formula} is \code{Expr} with: atoms instead of numbers, \code{¬} instead of \code{neg}, and \code{∧ ∨ →} instead of \code{plus}/\code{times}. \code{subformulas} is \code{subexprs}; \code{depth} is \code{height}. The work is in matching \emph{Mileti's} exact definitions, which is why Part~4 comes before Part~5.
\end{tip}

\subsection{Lesson F: notation, writing trees like math (30 min)}
Writing \code{plus (num 1) (times (num 2) (num 3))} gets tiring. Notation lets you write \code{num 1 ⊕ num 2 ⊗ num 3}.
\begin{lean}
-- (inside namespace Demo, after `open Expr`)
prefix:75 "−" => Expr.neg
infixl:65 " ⊕ " => Expr.plus
infixl:70 " ⊗ " => Expr.times      -- 70 > 65: ⊗ binds tighter than ⊕

#eval num 1 ⊕ num 2 ⊗ num 3
example : (num 1 ⊕ num 2 ⊗ num 3) = plus (num 1) (times (num 2) (num 3)) := by rfl
example : (−num 1 ⊗ num 2) = times (neg (num 1)) (num 2) := by rfl
\end{lean}
\begin{output}
Demo.Expr.plus (Demo.Expr.num 1) (Demo.Expr.times (Demo.Expr.num 2) (Demo.Expr.num 3))
\end{output}
Grouping to the right with \code{infixr}:
\begin{lean}
namespace R
inductive T where
  | leaf : Nat → T
  | arr  : T → T → T
  deriving Repr, DecidableEq
infixr:25 " ⇒ " => T.arr
open T
example : (leaf 1 ⇒ leaf 2 ⇒ leaf 3) = arr (leaf 1) (arr (leaf 2) (leaf 3)) := by rfl
end R
\end{lean}
\textbf{The rules:}
\begin{itemize}
\item \code{prefix:N "sym" => f}: you may write \code{sym x} for \code{f x}.
\item \code{infixl:N "sym" => f} / \code{infixr:N "sym" => f}: you may write \code{a sym b} for \code{f a b}.
\item \textbf{Precedence:} a higher \code{N} binds tighter. \code{⊗} at 70 beats \code{⊕} at 65, just as $\times$ beats $+$.
\item \textbf{Associativity:} with several in a row, \code{infixl} groups from the left, $(a\oplus b)\oplus c$, and \code{infixr} groups from the right, $a\Rightarrow(b\Rightarrow c)$.
\item \textbf{Test notation with \code{example … := by rfl}}: write the notation on the left and the constructor form on the right. If the test passes, the notation parses the way you meant.
\end{itemize}
\begin{warn}[Choose new symbols for your formulas]
Lean already uses \code{¬ ∧ ∨ →} for its \emph{own} logic. If you reuse them, your notation clashes with Lean's. Pick different symbols, for example \code{∼ ⋀ ⋁ ⟶} (tested; they work). If Lean says \code{unexpected token}, the symbol is taken or clashes; pick another.
\end{warn}

% ------------------------------------------------------------------
\section{Read Mileti with a purpose (45--60 minutes)}
Read \textbf{Ch.~2} (generating systems, especially ``freely generated'' and definitions by recursion) and \textbf{§3.1} (the definition of formulas). Keep a notes page and answer these questions. Write the answers in today's diary entry.

\begin{enumerate}[leftmargin=*]
\item \textbf{Which connectives are primitive in §3.1?} The project's design assumes $\neg,\wedge,\vee,\to$. Check it. Is $\leftrightarrow$ a primitive or an abbreviation? Is there a $\bot$ symbol? \emph{Only primitives become constructors.}
\item \textbf{What are the atoms?} Mileti uses a set $P$. Is $P$ fixed, countable, arbitrary? This affects Part~5, Step 1 (\code{Nat} vs.\ \code{String}).
\item \textbf{How does Mileti define the subformulas of a formula?} Write the definition and section number. Is the result a set?
\item \textbf{Does Mileti define depth or height (or a similar measure)?} Write the definition. Is an atom 0 or 1?
\item \textbf{In your own words: why is an inductive type a free generating system?} Use Part~1.2 of this guide. Write 4--6 sentences.
\end{enumerate}
\begin{tip}[Why question 5 matters]
It's ready-made material for your statement of purpose and interviews: ``Formalizing Mileti's syntax in Lean made me see why unique readability is needed for strings but automatic for inductive types…'' Write it carefully.
\end{tip}

% ------------------------------------------------------------------
\section{Build \texttt{Syntax.lean}, one small step at a time}
Open \code{lean/MiletiLogic/Syntax.lean}. It already contains a checklist comment and an empty namespace. Keep the checklist and tick items as you go (change \code{[ ]} to \code{[x]}). Write all code \textbf{between} \code{namespace MiletiLogic} and \code{end MiletiLogic}.

\begin{okbox}[After every step]
\begin{enumerate}
\item Save (⌘S).
\item The editor has no red underlines and the infoview shows no errors.
\item In Terminal: \code{cd \textasciitilde/Desktop/grad-school-project/mileti-kg-prover/lean \&\& ~/.elan/bin/lake build} prints \code{Build completed successfully}.
\item Commit (Part~6.2).
\end{enumerate}
\end{okbox}

\subsection{Step 1: decide how atoms are named (10 min)}
\begin{center}
\begin{tabularx}{\linewidth}{@{}l>{\raggedright\arraybackslash}X>{\raggedright\arraybackslash}X>{\raggedright\arraybackslash}X@{}}
\toprule
\textbf{Option} & \textbf{An atom looks like} & \textbf{Good} & \textbf{Less good}\\
\midrule
\code{Nat} & \code{atom 0}, \code{atom 1} & Simplest; easy to make many atoms; easy proofs & The Python harness writes \code{A}, \code{B}, so the export needs a mapping (A→0, B→1)\\
\code{String} & \code{atom "A"} & Matches the Python harness exactly & Slightly clunkier to type\\
\bottomrule
\end{tabularx}
\end{center}
Choose one. \textbf{Write a comment above your type explaining why.} The reason matters more than the choice, and you'll need it in your diary.

\subsection{Step 2: write \texttt{inductive Formula} (20--30 min)}
Model it on \code{Expr} from Lesson E:
\begin{itemize}
\item one constructor for \textbf{atoms}, carrying a \code{Nat} or \code{String} (like \code{num});
\item one for \textbf{¬}, carrying one \code{Formula} (like \code{neg});
\item one each for \textbf{∧}, \textbf{∨}, \textbf{→}, carrying two \code{Formula}s (like \code{plus});
\item \code{deriving DecidableEq, Repr};
\item a comment with the Mileti section it follows.
\end{itemize}
Name constructors clearly, for example \code{atom}, \code{not}, \code{and}, \code{or}, \code{imp}.
\begin{okbox}
\begin{enumerate}
\item With \code{\#eval}, build $(A\wedge B)\to(A\vee B)$ \emph{using constructors only} (no notation yet) and look at the printed tree. Does it have the shape of the tree in Part~1.2?
\item Add \code{example : <your first atom> ≠ <your second atom> := by decide}.
\item Add \code{example : <A ∧ B as constructors> ≠ <B ∧ A as constructors> := by decide}.
\end{enumerate}
\end{okbox}
\begin{tip}[Names]
Names like \code{not}, \code{and}, \code{or} are fine as constructor names (tested: they don't clash after \code{open Formula}).
\end{tip}
\textbf{Commit now.} This is your first real milestone.

\subsection{Step 3: notation (30--45 min)}
Following Lesson F:
\begin{itemize}
\item one symbol each for ¬, ∧, ∨, → (new symbols, not Lean's own);
\item precedences that match the Python harness, so the same text means the same formula on both sides:
\textbf{¬ binds tightest, then ∧, then ∨, then → loosest}, and \textbf{→ groups to the right} (\code{infixr}).
\end{itemize}
\begin{okbox}[Three precedence tests]
Write one \code{example … := by rfl} for each, with notation on the left and the constructor form on the right:
\begin{enumerate}
\item $\neg A\wedge B$ is $(\neg A)\wedge B$, not $\neg(A\wedge B)$.
\item $A\wedge B\vee C$ is $(A\wedge B)\vee C$.
\item $A\to B\to C$ is $A\to(B\to C)$.
\end{enumerate}
If a test fails, your precedence numbers or \code{infixl}/\code{infixr} choice is wrong. Good: the test caught it.
\end{okbox}

\subsection{Step 4: \texttt{subformulas} (20--30 min)}
Write \code{subformulas : Formula → List Formula}, modelled on \code{subexprs}, matching Mileti's definition (Part~4, question 3).
\begin{okbox}
Take $\varphi = \neg(A\wedge B)\to A$.
\begin{enumerate}
\item On paper, list its subformulas in the order your function will produce them.
\item \code{\#eval subformulas φ} and compare.
\item Count them; add \code{example : (subformulas φ).length = <your count> := by rfl}.
\end{enumerate}
\end{okbox}
\textbf{Decision to record:} $A$ appears twice in $\varphi$. Does your list keep duplicates? Does Mileti's definition give a set? Write your choice and reason in a comment. (\code{.eraseDups} gives set-like behaviour if you want it.)

\subsection{Step 5: \texttt{depth} (15--20 min)}
Write \code{depth : Formula → Nat}, modelled on \code{height}, matching Mileti's definition (Part~4, question 4). If the book doesn't define one, use height with atoms at 0 and say so in a comment.
\begin{okbox}
Compute by hand, then add an \code{example … := by rfl} for: an atom; $\neg A$; $(A\wedge B)\to(A\vee B)$; $\neg\neg\neg A$.
\end{okbox}

\subsection{Step 6: clean up (15 min)}
\begin{itemize}
\item The comment at the top says what's in the file, with Mileti section references, and the checklist is ticked.
\item Each decision (atoms, symbols, list vs.\ set, depth convention) has a one-line comment with its reason.
\item There's no \code{sorry} anywhere. (\code{sorry} means ``unfinished'' to Lean.)
\item \code{lake build} passes.
\end{itemize}

% ------------------------------------------------------------------
\section{Finish: checklist and git}

\subsection{``Done'' checklist for step 1}
\begin{itemize}[label=□]
\item \code{Formula} with atoms and ¬ ∧ ∨ →, \code{deriving DecidableEq, Repr}
\item Atom choice explained in a comment
\item Notation with three passing precedence tests
\item \code{subformulas}, checked by hand on one formula, duplicates decision recorded
\item \code{depth}, checked by hand on four formulas
\item At least two \code{≠ … := by decide} tests
\item \code{lake build} passes, no \code{sorry}
\item Committed and pushed; diary entries for each working day
\end{itemize}

\subsection{Committing and pushing}
A \textbf{commit} saves a snapshot of your work in git's history; \textbf{push} uploads it to GitHub. From the project folder:
\begin{shell}
cd ~/Desktop/grad-school-project/mileti-kg-prover
git status                                   # what changed? (red = not staged)
git add lean/MiletiLogic/Syntax.lean         # stage your file
git commit -m "Syntax: Formula type with atoms and ¬ ∧ ∨ →"
git push                                     # upload to GitHub
\end{shell}
\begin{tip}[Good commit messages]
Say \emph{what} changed: ``Syntax: notation with precedence tests'', ``Syntax: subformulas and depth''. Commit after each step in Part~5. Many small commits are visible evidence of steady, independent work, and admissions readers can see them on GitHub.
\end{tip}
\begin{okbox}
Open \code{github.com/ZaDrMeister/mileti-kg-prover} in a browser; your commit message appears at the top of the file list.
\end{okbox}

% ------------------------------------------------------------------
\section{When something goes wrong}

\subsection{Real Lean error messages}
\begin{center}
\begin{tabularx}{\linewidth}{@{}>{\raggedright\arraybackslash}p{4.3cm}>{\raggedright\arraybackslash}X>{\raggedright\arraybackslash}X@{}}
\toprule
\textbf{You see} & \textbf{It means} & \textbf{Fix}\\
\midrule
\code{Missing cases: …} & A function doesn't handle every constructor & Add a line for each constructor listed\\
\code{failed to synthesize Decidable (…)} & Lean can't decide this equality & Add \code{deriving DecidableEq} to the type\\
\code{Unknown identifier `atom`} & The name isn't in scope & Write \code{Formula.atom}, or add \code{open Formula}\\
\code{unexpected token} at your symbol & The symbol is already used or clashes & Choose another symbol\\
\code{Tactic `rfl` failed: The left-hand side … is not definitionally equal to the right-hand side} & The two sides are different trees & Check precedence, or redo your hand calculation\\
Infoview empty, nothing happens & Lean is starting, or you opened the wrong folder & Wait 30 s; make sure you opened the \code{lean} folder\\
\bottomrule
\end{tabularx}
\end{center}

\subsection{How to ask for help (from Claude or a person)}
Paste three things:
\begin{enumerate}
\item the code (the smallest piece that shows the problem);
\item the \emph{full} error message;
\item what you expected to happen.
\end{enumerate}
Then ask ``What does this error mean?'' or ``Give me a hint.'' Not ``fix it''. You need to be able to explain every line of this file.

\subsection{References}
\begin{itemize}
\item \emph{Functional Programming in Lean} (lean-lang.org): the most beginner-friendly; read the chapters on structures, datatypes and pattern matching.
\item \emph{Theorem Proving in Lean 4} (lean-lang.org): chapters ``Inductive Types'' and ``Induction and Recursion''.
\item Avigad, \emph{Logic and Proof}, ch.~4 ``Propositional Logic in Lean'': connects natural deduction and Lean.
\item For step 2: \emph{Intuitionistic Propositional Logic in Lean} (arXiv 2410.23765) encodes a proof system as an inductive type with one constructor per rule, the same pattern you'll use.
\end{itemize}

% ------------------------------------------------------------------
\section{Your research diary}
At the end of every working day:
\begin{shell}
cd ~/Desktop/grad-school-project/mileti-kg-prover
./diary/new-entry.sh          # creates diary/entries/YYYY-MM-DD.md
\end{shell}
Open the file and fill in: goal, what you did, evidence (commits, outputs), what you learned, \textbf{decisions and why}, open questions (open questions marked with a question mark, resolved ones with ✓), AI help used, and the next step. Then commit it:
\begin{shell}
git add diary && git commit -m "Diary: 2026-11-26" && git push
\end{shell}
\textbf{What to write for step 1:}
\begin{itemize}
\item Day 1: setup problems and how you solved them; anything in Lessons A--C that surprised you.
\item Day 2: your answers to the five Mileti questions (Part~4), especially question 5.
\item Day 3: your decisions (atoms, symbols, list vs.\ set, depth) with reasons; the commits you made.
\end{itemize}

% ------------------------------------------------------------------
\section{Self-test and next steps}

\subsection{Can you explain these without notes?}
If you can, you're ready for step 2. They're also likely interview questions.
\begin{enumerate}
\item What does \code{inductive} do, and why does it give unique readability for free?
\item Why does Lean accept \code{size} calling itself?
\item What do \code{deriving DecidableEq} and \code{deriving Repr} each do, and why does the project need \code{DecidableEq}?
\item What's the difference between \code{by rfl} and \code{by decide}?
\item Why did you choose your precedences, and how do you know they're right?
\item Are your subformulas a list or a set, and why?
\end{enumerate}

\subsection{Step 2 preview: \texttt{Deduction.lean}}
Next you'll write
\begin{center}\code{inductive Derivable : List Formula → Formula → Prop}\end{center}
with \textbf{one constructor per rule} from \code{docs/design.md} §3. It's the same \code{inductive} idea, but instead of describing what a formula \emph{is}, the type describes what can be \emph{derived}: \code{Derivable S φ} is the Lean version of ``$S\vdash\varphi$ has a deduction.'' You'll start with Assume, ∧EL, ∧ER, ∧I, →E and →I, and prove the first book item, $A\wedge B\vdash A\vee B$ (which also needs ∨IL). A separate guide will come when step 1 is done.

\end{document}
